\section{Specifika synchronizace dat}
\label{sec:implementation:gotchas}

Tato sekce popisuje konkrétní implementační přístupy a úskalí při synchronizaci dat z~GitHubu a~z~lokálního git repozitáře.

\subsection{Inkrementální načítání}
\label{subsec:implementation:gotchas:incremental}

Při prvním spuštění synchronizace je nutné načíst historická data. Repozitář \code{rust-lang/rust} v~době psaní této práce obsahuje přibližně 92~000 \PR a~61~000 issues. Opakovat takový plný průchod při každém obnovení dat by bylo jak časově nákladné, tak zbytečně zatěžující z~pohledu počtu požadavků na \API GitHubu.

Z~tohoto důvodu systém rozlišuje dva režimy synchronizace (\code{SyncMode}):

\begin{itemize}
  \item \textbf{Last(N)} — načte posledních $N$ stránek výsledků (\PR nebo issues), seřazených sestupně podle data vytvoření. Slouží pro rychlé doplnění nejnovějších záznamů nebo při prvotním spuštění.
  \item \textbf{Since(datum)} — načítá záznamy seřazené sestupně podle časového razítka a~ukončí se ve chvíli, kdy narazí na záznamy starší než zadané datum. Tento režim je využíván při průběžné synchronizaci v~pozadí.
\end{itemize}

V~rámci serveru běží na pozadí smyčka (\code{StateMonitor}), která se každých 60 sekund spustí a~zavolá synchronizaci v~režimu \code{Since(last\_known\_event\_timestamp)}, kde jako počáteční datum slouží časová značka posledního záznamu v~tabulce \code{issue\_event\_history}. Tím je zajištěno, že jsou průběžně doplňovány pouze záznamy, které se od posledního spuštění změnily.

Pro případ, že od počátečního naplnění databáze chybí historické záznamy (například protože při prvním synchronizaci nebyla volána možnost \code{--full-history}), je k~dispozici příkaz \code{backfill}. Ten vyhledá záznamy v~tabulce \code{issues}, ke kterým dosud v~tabulce \code{issue\_event\_history} neexistují žádné řádky, a~pro každý takový záznam zpětně načte celou historii událostí z~GitHub \API.

\subsection{Limity GitHub \API a jejich vliv na návrh}
\label{subsec:implementation:gotchas:rate_limits}

GitHub \REST \API pro autentizované uživatele umožňuje 5~000 požadavků za hodinu\cite{github_rest_api}. V~kontextu průběžné synchronizace, kdy jsou každých 60 sekund zasílány pouze požadavky pro záznamy s~nedávnou aktivitou, je reálné zatížení výrazně pod tímto limitem, a~to v~řádu desítek až stovek požadavků za hodinu. Limity tedy nejsou při normálním provozu překračovány.

Přesto je v~implementaci přítomna ochrana pro případ rozsáhlé počáteční synchronizace nebo zpětného doplnění (\code{backfill}), kdy se může počet požadavků přiblížit limitům. Pro každý požadavek na \API GitHubu je použita funkce \code{retry\_on\_rate\_limit}, která detekuje chyby odpovídající překročení limitu (HTTP status 429, řetězce \code{"rate limit"}, \code{"too many requests"} nebo \code{"secondary rate"} v~chybové zprávě)\cite{github_rest_api} a~automaticky opakuje požadavek s~exponenciálním zpožděním. Počáteční čekání je 30 sekund, maximální prodleva je 15 minut a~celková maximální doba opakování je 45 minut. Pokud se požadavek nezdaří ani po uplynutí této doby, chyba se propaguje dál a~synchronizace je přerušena. Tento mechanismus slouží jako pojistka, nikoli jako primární strategie pro správu limitu.

\subsection{Timeline events vs.\ webhooky}
\label{subsec:implementation:gotchas:timeline_vs_webhooks}

GitHub nabízí dva přístupy k~získání průběžné historie změn \PR a~issues: webhooky a~timeline events \API.

Webhooky fungují jako push notifikace — GitHub při každé relevantní události zašle \HTTP požadavek na předem registrovaný endpoint. Tento přístup je vhodný pro aplikace, které potřebují reagovat na události v~reálném čase\cite{github_webhooks}. Pro účely této práce je však nevhodný ze dvou důvodů. Za prvé, webhook vyžaduje veřejně dostupný endpoint, což komplikuje lokální vývoj a~testování. Za druhé, webhooky neposkytují historická data, takže jimi nelze zpětně získat události, ke kterým došlo před registrací webhooku.

Systém proto používá timeline events \API, které pro každý \PR nebo issue vrátí seznam všech zaznamenaných událostí v~chronologickém pořadí. Tím je možné jak zpětně naplnit historii, tak průběžně doplňovat nové události. Načítání je optimalizováno tak, že pro každý záznam je kontrolována časová značka posledního uloženého záznamu v~databázi.  Pokud načtená stránka událostí obsahuje záznamy s~časovou značkou starší než tato hodnota, načítání je ukončeno.

\subsection{Kombinace GitHub \API a~git repozitáře}
\label{subsec:implementation:gotchas:api_vs_git}

Informace o~souborech změněných v~rámci sloučeného \PR jsou dostupné přes \API GitHubu prostřednictvím endpointů pro výpis souborů daného \PR\cite{github_rest_api}. Použití tohoto endpointu by však při zpracování tisíců \PR přineslo výrazné navýšení počtu požadavků na \API, a taktéž by bylo pomalé kvůli síťové latenci.

Z~tohoto důvodu jsou informace o~změněných souborech získávány přímo z~lokální kopie git repozitáře pomocí knihovny \code{git2}. Pro každý sloučený \PR je znám \SHA merge commitu, který je uložen v~tabulce \code{issues}. Pomocí funkce \code{diff} je tento commit porovnán s~jeho prvním rodičovským commitem, čímž vzniká seznam souborů změněných v~rámci daného \PR. Tento přístup nevyžaduje žádné síťové požadavky na \API GitHubu a~je výrazně rychlejší, protože git repozitář je k~dispozici lokálně.

Nevýhodou lokálního přístupu je nutnost mít k~dispozici naklonovaný repozitář \code{rust-lang/rust}, který má v~době psaní této práce velikost přibližně 6~GB (samotná git historie bez zdrojových souborů). Systém repozitář při prvním spuštění automaticky naklonuje a~při každé synchronizaci stáhne nejnovější změny pomocí \code{git fetch}.
